\documentclass[amsmath, amssymb, aps]{revtex4-2}

\usepackage{graphicx}
\usepackage{dcolumn}
\usepackage{bm}
\usepackage{booktabs}
\usepackage{hyperref}
\usepackage{algorithm}
\usepackage{algpseudocode}
\usepackage{caption} 
\captionsetup[algorithm]{labelfont=bf, textfont=bf}
\usepackage{float}
\usepackage{placeins}
\usepackage{geometry}
\usepackage{setspace}
\usepackage{etoolbox}
\makeatletter
\@removefromreset{equation}{section}
\makeatother
\geometry{margin=1in}
\onehalfspacing

\hypersetup{
    colorlinks=true,
    linkcolor=blue,
    filecolor=magenta,      
    urlcolor=cyan,
}


\setlength{\parskip}{1em}

\begin{document}

\pagestyle{plain}

\title{Four New Multi-Phase Hybrid Bracketing Algorithms for Numerical Root Finding }

\author{Abdelrahman Ellithy$^{1}$}
\affiliation{Department of Artificial Intelligence, Faculty of Computers and Artificial Intelligence, Benha University, Benha, Egypt}

\date{\today}

\begin{abstract}
Finding roots of nonlinear equations requires balancing rapid convergence with guaranteed bracketing reliability. This work introduces four innovative high-performance hybrid algorithms — \textbf{Opt.BF, Opt.BFMS, Opt.TF, and Opt.TFMS} — These methods systematically combine robust interval contraction with accelerated interpolation and modified-secant corrections. Evaluated on a diverse \textbf{48-problem} benchmark, the leading method, Opt.BFMS, converges in only 2.79 iterations on average (23 wins in fewest iterations) using 3–4 function evaluations per iteration (NFE/iter) and 48–58 FLOPs/iter excluding function evaluations, yielding the lowest total estimated workload (approximately \textbf{401–650 FLOPs} per problem). This translates to an average CPU time of 10.48 microseconds — a 3.5x speedup over Bisection, 1.91x over Brent (SciPy), and \textbf{2.54x} over QIR — while winning 45.83\% of problems by execution time and \textbf{47.9\% by total NFE and Iteration count as well}. These results establish the proposed hybrids, particularly Opt.BFMS, and Opt.TFMS as the new state-of-the-art in practical root-finding performance across polynomial, transcendental, and ill-conditioned equations.
\end{abstract}


\maketitle

\noindent\textbf{Keywords:} hybrid root-finding algorithms, root finding, Continuous optimization, convergence analysis, Derivative-free methods, Bracketing methods, modified secant method, Nonlinear equations, numerical analysis, numerical stability.

\section{Declarations}\label{sec:declaration}

\subsection{Data availability}\label{sec:dataavailability}
The datasets generated and/or analyzed during the current study are available from the corresponding author on reasonable request. The SymPy implementations of all algorithms and experimental scripts are available at: \url{https://github.com/Abdelrhman-Ellithy/NM-Algorithms/}

\end{document} 